\begin{explanationbox}
	\textbf{\textcolor{CExplanation}{一句话直觉}}

	将点到平面的距离看作向量 $\vec{AP}$ 在平面法向量 $\mathbf{n}$ 方向上的投影长度，并取绝对值即可得到距离。

	\medskip
	\textbf{\textcolor{CExplanation}{核心拆解}}
	\begin{description}[style=nextline, leftmargin=2.8em, labelsep=0.5em, itemsep=0.45em]
		\item[\textbf{要点一（向量投影原理）：}] 空间任意直线（通过法向量方向）上的投影长度公式为 $\frac{|\vec{AP}\cdot\mathbf{n}|}{|\mathbf{n}|}$。
		\item[\textbf{要点二（绝对值必要性）：}] 投影方向可能与法向量同向或反向，距离是非负数，因此必须取绝对值。
	\end{description}

	\medskip
	\textbf{\textcolor{CExplanation}{几何本质（垂线段长度）}}

	点 $P$ 到平面 $\alpha$ 的最短距离就是过 $P$ 向平面作垂线，垂足与 $P$ 的连线长度。而法向量 $\mathbf{n}$ 恰好指示了这条垂线的方向，因此 $\vec{AP}$ 在 $\mathbf{n}$ 上的投影长度就是垂线段长度。

	\medskip
	\textbf{\textcolor{CExplanation}{代数意义（数量积运算）}}

	利用向量点积，将几何量（垂线段）转化为代数式 $\vec{AP}\cdot\mathbf{n}$，并通过除以 $|\mathbf{n}|$ 实现单位化，最终结果只含加减与绝对值运算，方便计算。

	\medskip
	\textbf{\textcolor{CExplanation}{考点价值（直接应用与逆向）}}
	\begin{itemize}[leftmargin=*, itemsep=0.5em, topsep=0.4em]
		\item \textbf{考法一（代入公式求距离）：} 给出平面的法向量或平面方程，以及平面上的点和外部点，直接套公式计算。
		\item \textbf{考法二（已知距离反求参数）：} 给出距离条件，利用公式构造方程，求解未知点坐标或法向量参数。
	\end{itemize}

	\medskip
	\textbf{\textcolor{CExplanation}{顿悟点}}

	为什么用 $\vec{AP}$ 而不是点 $P$ 到平面上任意点的向量？因为距离与所取点 $A$ 无关（证明中需说明投影是定值），理解这一点才能灵活选点。

	\medskip
	\textbf{\textcolor{CExplanation}{使用场景}}
	\begin{itemize}[leftmargin=*, itemsep=0.5em, topsep=0.4em]
		\item \textbf{场景一（建系求点面距）：} 在立体几何题中，当能方便建立空间直角坐标系并求出法向量时，优先用此公式。
		\item \textbf{场景二（平行平面间距离）：} 将平行平面间的距离转化为一个平面内一点到另一平面的距离，同样适用此公式。
	\end{itemize}
\end{explanationbox}
